% Section: proven-optimal references (Yoon).
% Needs \usepackage{booktabs, amsmath}. Tables are generated by report/make_tables.py.

\subsection{Proven-optimal reference plans}
\label{sec:exact}

To say how far an LLM's plan is from the best possible plan we need the optimum of every
instance, not just a good plan. A reference that is itself sub-optimal makes the comparison
misleading: against our first references, the direct plans of each Claude model appeared to \emph{beat} the
reference on 20 of the 22 field reports for the 11 instances that were not yet proven
optimal, which says more about the reference than about the LLM.

\paragraph{Fast Downward alone is not enough.}
We first solved every instance with Fast Downward~\cite{helmert2006fd}, using the optimal
configurations A$^*$ with LM-cut~\cite{helmert2009lmcut} and A$^*$ with merge-and-shrink
(M\&S)~\cite{helmert2014ms}, 300\,s each, falling back to anytime LAMA~\cite{richter2010lama}
when both ran out of time. Table~\ref{tab:proof-summary} shows that this proves every small
instance and 14 of 15 medium ones optimal, but only 5 of 15 large ones: 34 of 45 overall.
For the 11 unproven instances LAMA's best plan was 17--88\% above the optimum
(Table~\ref{tab:unproven}). Taking the cheaper of the splitting and no-split domains, which was
our reference before the exact solver, narrowed this to 0--45\%, but those references still
could not be shown to be optimal.

The difficulty comes from the encoding rather than from the size of the road network.
Counts of people are level objects with a static \texttt{(sum a b c)} relation, so
\texttt{load} has ten parameters and grounds to many actions, and the same set of rescues can
be reached through many orders of loads, drives and car assignments. The heuristics are also
weak on this domain: under the delete relaxation a car keeps every \texttt{free-seats} value it
has ever had, so seat capacity disappears and one relaxed tour can pick up everybody. LM-cut
therefore ignores the repeated trips that make up most of the real cost on large instances.
Merge-and-shrink expanded up to 29.5 million states (medium-04) before proving optimality, and
LAMA expanded up to 3.3 million (large-15) without reaching the optimum.

\paragraph{An exact domain-specific solver.}
We therefore wrote \texttt{exact.py}, a uniform-cost (Dijkstra) search over a much smaller
state space. Write $G=(V,E)$ for the graph of undamaged roads with travel times $t(e)$,
$d(u,v)$ for shortest-path distances (Floyd--Warshall), $Q=4$ for seats per car and $p_b$ for
the deferral penalty of band $b$. The objective is the total driving time plus
$\sum p_b$ over everyone deferred. Five properties of the domain let us shrink the search
without losing optimality:
\begin{enumerate}
  \item \emph{Shortest paths.} Between two stops (a pick-up or a drop-off) a car should
    follow a shortest path, so driving is precomputed as $d(u,v)$.
  \item \emph{Trips.} Unloading is free, only possible at evacuation points, and every car
    must end empty, so each car's route is a sequence of trips: from its current position,
    visit a set $S$ of pick-up sites, then drive to an evacuation point $e$. A trip costs the
    shortest route through $S$ to $e$, found by trying every order ($|S| \le Q$, so at most
    24 orders).
  \item \emph{Cars one at a time.} The model has no clock, the cost is a plain sum, and cars
    only interact through who is still waiting. Any joint plan can be reordered so car 1
    does all its trips, then car 2, at the same cost.
  \item \emph{Most urgent first.} At one site, taking a higher-penalty person instead of a
    lower-penalty one never raises the cost, so we only need to know \emph{how many} people
    $n_v$ have been taken from site $v$; they are taken in decreasing order of $p_b$.
  \item \emph{Full cars.} Taking one more person on a trip that already visits their site
    adds no driving and removes a penalty, so a trip through $S$ carries
    $\min(Q, \sum_{v\in S} \text{waiting}_v)$ people.
\end{enumerate}
A state is then $(c, \mathit{pos}, (n_v)_{v})$: the car being planned, where it is, and how
many people have been taken from each site. The successors are every trip of car $c$
(cost: its driving), handing over to car $c{+}1$ (cost 0), and, after the last car,
deferring everyone left (cost: their penalties). All edge costs are non-negative, so Dijkstra
returns an optimal plan. We do not need a heuristic: the reduced space has at most 242{,}399
states (large-03), against the millions Fast Downward expanded.

\begin{table}[t]
  \centering\small
  \begin{tabular}{lrrrrr}
\toprule
 & \multicolumn{2}{c}{Fast Downward} & \multicolumn{3}{c}{\texttt{exact.py}} \\
\cmidrule(lr){2-3}\cmidrule(lr){4-6}
Size & proven & by LM-cut / M\&S & proven & max time (s) & max states \\
\midrule
Small & 15/15 & 15 / 0 & 15/15 & $<$0.1 & 47 \\
Medium & 14/15 & 10 / 4 & 15/15 & 0.4 & 5{,}197 \\
Large & 5/15 & 2 / 3 & 15/15 & 109.0 & 242{,}399 \\
\midrule
All & 34/45 & & 45/45 & & \\
\bottomrule
\end{tabular}

  \caption{Proving the 45 reference plans optimal. Fast Downward ran A$^*$ with LM-cut, then
  A$^*$ with merge-and-shrink, with a 300\,s limit for each search. \texttt{exact.py} times are
  single-threaded Python on one core.}
  \label{tab:proof-summary}
\end{table}

\begin{table}[t]
  \centering\small
  \begin{tabular}{lrrrrr}
\toprule
Instance & Optimum & FD split (LAMA) & gap & Old reference & gap \\
\midrule
medium-14 & 401 & 471 & 17\% & 401 & 0\% \\
large-01 & 240 & 452 & 88\% & 268 & 12\% \\
large-02 & 428 & 721 & 68\% & 461 & 8\% \\
large-03 & 434 & 602 & 39\% & 542 & 25\% \\
large-06 & 445 & 646 & 45\% & 446 & 0\% \\
large-08 & 366 & 575 & 57\% & 430 & 17\% \\
large-10 & 402 & 616 & 53\% & 480 & 19\% \\
large-11 & 403 & 586 & 45\% & 529 & 31\% \\
large-12 & 302 & 437 & 45\% & 437 & 45\% \\
large-13 & 212 & 389 & 83\% & 250 & 18\% \\
large-15 & 476 & 695 & 46\% & 487 & 2\% \\
\midrule
Mean & & & 53\% & & 16\% \\
\bottomrule
\end{tabular}

  \caption{The 11 instances Fast Downward could not prove optimal in 300\,s, with the optimum
  found by \texttt{exact.py}. ``Old reference'' is the cheaper of the Fast Downward plans for
  the splitting and no-split domains, which we used before the exact solver.}
  \label{tab:unproven}
\end{table}

\paragraph{Checking the solver.}
Because \texttt{exact.py} is our own code, we checked it in two independent ways. It returns
exactly the same cost as Fast Downward on all 34 instances that Fast Downward proved optimal,
and every plan it outputs is accepted by our plan checker (which agrees with VAL) with the
claimed cost. All 45 instances solve in about four minutes in total, most in under a second
(Table~\ref{tab:proof-summary}); the per-instance optima are listed in
Appendix~\ref{app:optima}. These 45 proven optima are the references used throughout our
evaluation.

\paragraph{What this says about the tool.}
The exact solver is not a replacement for a planner in our comparison: it is hand-written for
this one domain, and the reductions above stop being exact as soon as the model gains a clock
(deadlines, or cars that must coordinate in time), non-identical seats, or a cost to unload.
It does show that the hard part of these instances is search, not modelling: a general
planner given a correct PDDL model still fails to prove optimality on 11 of 45 instances
within five minutes, while a solver that exploits the structure of the problem needs seconds.
This is the cost of generality the PDDL pipeline pays, and it matters when we compare it with
the LLM, which plans directly and can use the same kind of structure implicitly.
